\section{Introduction}\label{sec:intro}

Referees are agents whose decisions are subjective, consequential, and observed by large and partisan audiences. A large literature documents that these decisions respond to social pressure, most prominently in the form of home bias in stoppage time, cards, and penalties \citep{sutter2004favoritism, garicano2005favoritism, boyko2007referee, dohmen2008influence, rickman2008favouritism, scoppa2008subjective, buraimo2010twelfth}. Natural experiments that removed crowds, from stadium closures in Italy to the Covid-19 ``ghost games,'' show that much of this bias disappears without spectators \citep{pettersson2010behavior, endrich2020home, bryson2021causal, reade2022eliminating}. Explicit corruption has also been documented: the 2006 \textit{Calciopoli} scandal in Italy revealed referee designations steered toward favored clubs \citep{boeri2011match}.

Much less is known about whether particular \emph{clubs}, as opposed to home teams, receive favorable treatment, and how large such an advantage is relative to every other club. The question is pressing in Spain. In 2023, prosecutors alleged that FC Barcelona had paid about XX [verify: reported as roughly \texteuro 7 million] between 2001 and 2018 to companies owned by José María Enríquez Negreira, then vice-president of the Comité Técnico de Árbitros (CTA), the body that oversees Spanish referees.\footnote{XX: Add the source for the prosecutor's complaint (March 2023) and the reported amounts. Whether any refereeing decision was affected is the subject of ongoing legal proceedings; this paper makes no claim about the legal case.} Whether those payments bought anything on the pitch is ultimately an empirical question about refereeing decisions.

In this paper, I ask three questions. First, do referees in the top-5 European leagues make systematically more favorable decisions for some clubs than for others? Second, where does Barcelona sit in that distribution, compared with Real Madrid, the other elite clubs, and every other club? Third, is Barcelona treated differently by Spanish referees than by UEFA referees, and did that difference change after 2017/18?

To answer these questions, I assemble match-level data on fouls, cards, penalties, possession, and stoppage time for about 48,500 matches in the Premier League, La Liga, Serie A, the Bundesliga, Ligue 1, and the Champions League between 2001/02 and 2025/26. I link these data to pre-match bookmaker odds. I estimate club-specific decision gaps conditional on opponent, home status, league-season, and pre-match win probability. I place Barcelona's estimate in the distribution of estimates for all 144 clubs with at least 150 matches, which serves as a permutation-style benchmark. To separate club-level playing style from referee behavior, I exploit the fact that Barcelona's style travels to the Champions League while its Spanish referees do not. I also normalize fouls by possession, which measures a side's exposure to being fouled.

Preview of findings. Raw comparisons strongly favor Barcelona, but they mostly reflect its possession-dominant style. Conditional on fouls called, Barcelona's players are no less likely to be booked, and its opponents are no more likely. On penalties, Barcelona ranks 93rd of 150 clubs. Normalizing fouls by possession reverses the sign of its foul advantage. Barcelona's net yellow-card advantage in La Liga does decline after 2017/18 relative to Real Madrid, by 0.57 cards per match. That decline coincides with Barcelona's sporting and financial crisis, and its net-foul advantage breaks only in 2021/22, the season Messi left. The clearest evidence of favoritism is in stoppage time: referees add about 0.7 more minutes when Barcelona trails by a goal than when it leads, but Real Madrid enjoys exactly the same advantage.

This paper contributes to the literature on favoritism and social pressure in refereeing by moving from home bias to club-specific bias, and by benchmarking a single club against the full cross-section of clubs. It also contributes to the literature on corruption in sports \citep{boeri2011match} by showing how far observable decisions can and cannot speak to an alleged corrupt relationship.

The rest of the paper is organized as follows. Section \ref{sec:data} describes the data. Section \ref{sec:strategy} presents the empirical strategy. Section \ref{sec:results} reports the results. Section \ref{sec:conclusion} concludes.
